\section{Implementation}
\label{sec:implementation}

We implemented QPASE in Python~\MachinePython{} with numpy for the ring arithmetic, \texttt{cryptography} for
AES-256-GCM and \texttt{dilithium-py} for ML-DSA-44. The code in \texttt{src/ris\_sem\_project/} covers Setup, the
key generation, Register, Login, Outsource, single-keyword Retrieve and the key update. The \texttt{primitives/}
package holds the ring, sampling, the PRF, the TOPRF, HKDF, the extractor and the wrappers; \texttt{parties/} holds
the user and server classes and the wire messages; and \texttt{protocol.py} runs each phase and records every
message. The notebooks call this library and add no protocol code.

\paragraph{Exact ring arithmetic.} The modulus $q = 2^{28} - 57$ does not suit the number-theoretic transform, so we
multiply in $\Rq$ with a twisted complex FFT. Each operand is split into two 14-bit limbs, which keeps every
floating-point product exact. This works for any $n$, including the odd $n = 595$. Notebook~01 finds
\RingProductsExact{} random products equal to an exact reference.

\subsection{Where the paper is incomplete}
\label{sec:decisions}

The paper leaves several values open and contradicts itself in a few places. \Cref{tab:decisions} lists the gaps
that matter most. \texttt{docs/implementation-plan.md} gives the full list and the reasoning.

\begin{table}[ht]
\centering
\small
\caption{Main gaps in the paper and our choices.}
\label{tab:decisions}
\begin{tabularx}{\linewidth}{>{\raggedright\arraybackslash}p{0.16\linewidth}>{\raggedright\arraybackslash}X>{\raggedright\arraybackslash}X}
\toprule
\textbf{Topic} & \textbf{What the paper says} & \textbf{Our choice} \\
\midrule
Ring or matrices & Fig.~1 uses LWE matrices; Sections II-B and IV-A use vectors in $\Rq^{1\times\ell}$. & The
ring $\Rq = \mathbb{Z}_q[X]/(X^n+1)$. \\
TOPRF sign & Theorem~1 subtracts $2r\sum_i \lambda_i e_i$ but defines $e''$ with $+r\sum_i \lambda_i e_i$. & The
sign the algebra gives, \cref{eq:toprf-combine}. \\
Lagrange coefficients & Not discussed; $e''$ is small only if every $\lambda_i$ is an integer. & Only subsets with
integer coefficients; $\{1, 3\}$ with $t = 2$ is refused. \\
Extractor hint & Fig.~2 computes the hint from $vk_i$, unrelated to $F$. & $\sigma = S(F)$ at Register, stored by the
servers and returned at Login. \\
Extractor formula & Definition~8 prints $E(x, \sigma) = (x + \sigma)(q-1)/2 \bmod 2$. & Ding et al.'s form, for which
robustness holds. \\
Sharing degree & Degree $t$, so $t + 1$ servers reconstruct, but Login contacts $t$. & Degree $t - 1$. \\
Open values & $p$, $\sigma$, the encryption mode, the Dilithium level and the password encoding are not fixed. &
$p = \ParamP$, $\sigma = \ParamSigma$, $L = \ParamL$ bits of SHA3-256 of the password, AES-256-GCM, ML-DSA-44. \\
\bottomrule
\end{tabularx}
\end{table}

\paragraph{How we checked it.} \UnitTests{} unit tests check observable behaviour: exact ring products, the same
$\Ku$ across logins and server subsets, rejected wrong passwords and tampered messages, and documents recovered
byte for byte. Eight notebooks run each phase with real data and end with assertion cells. Every number in this
report comes from one run of all notebooks on one machine and is generated from the cell that printed it.
